% Title page and document layout follow the supplied report template.
% The title is inline so the open document is self-contained.
\documentclass[a4paper,12pt]{article}
\usepackage{iftex}
\ifPDFTeX
  \usepackage[utf8]{inputenc}
  \usepackage[T2A]{fontenc}
\else
  \usepackage{fontspec}
  \setmainfont{Liberation Serif}
  \setsansfont{Liberation Sans}
  \setmonofont{Liberation Mono}
\fi
\usepackage[russian]{babel}
\usepackage{geometry}
\usepackage{booktabs}
\usepackage{graphicx}
\usepackage{array,tabularx,float,xcolor}
\usepackage{tikz}
\usetikzlibrary{arrows.meta}
\geometry{left=2cm,right=2cm,top=2cm,bottom=2cm}

\newcommand{\reportcodesize}{\footnotesize}
\ifPDFTeX
  \usepackage{amsmath,amsfonts,amssymb}
  \usepackage{caption}
  \usepackage{hyperref}
  \hypersetup{colorlinks=true,linkcolor=black,urlcolor=blue}
  \usepackage{minted}
  \usemintedstyle{colorful}
  \setminted{fontsize=\reportcodesize,linenos,numbersep=7pt,
    breaklines,frame=single,framesep=2mm}
\else
  % The desktop editor compiles without external shell commands.
  % Keep the minted blocks compatible with that compiler as well.
  \usepackage{listings}
  \lstnewenvironment{minted}[2][]{\lstset{language=C,
    basicstyle=\ttfamily\reportcodesize,columns=fullflexible,
    keepspaces=true,showstringspaces=false,breaklines=true,
    frame=single,rulecolor=\color{black!35},numbers=left,
    numberstyle=\tiny\color{black!60},numbersep=7pt,
    xleftmargin=10pt,framexleftmargin=3pt,tabsize=2,
    keywordstyle=\bfseries\color[rgb]{0,0.53,0},
    commentstyle=\itshape\color{black!55},
    stringstyle=\color[rgb]{0.75,0.20,0},
    aboveskip=4pt,belowskip=6pt}}{}
\fi
\setlength{\parindent}{1.25cm}
\setlength{\parskip}{2pt}
\setlength{\emergencystretch}{2em}
\newcolumntype{Y}{>{\raggedright\arraybackslash}X}
\renewcommand{\arraystretch}{1.10}
\newcounter{reportlisting}
\newcommand{\codecaption}[1]{\refstepcounter{reportlisting}%
  \begin{center}\small Листинг \thereportlisting: #1\end{center}}
\newcommand{\code}[1]{\texttt{\detokenize{#1}}}
\tikzset{box/.style={draw,rounded corners=2pt,align=center,minimum height=10mm,
 text width=31mm,font=\small},arr/.style={-{Latex[length=2mm]},semithick},
 state/.style={draw,rounded corners=3pt,minimum width=28mm,minimum height=10mm,
 align=center,font=\small}}
\begin{document}
\begin{titlepage}
    \centering
    Федеральное государственное автономное образовательное учреждение высшего образования «Национальный исследовательский университет ИТМО» \\
    \vspace{1cm}
    Факультет Программной Инженерии и Компьютерной Техники \\
    \vfill
    {\large «Встроенные системы»} \\
    \vspace{0.5cm}
    Лабораторная работа №3 \\
    Программирование микроконтроллеров с операционной \\
    системой реального времени (FreeRTOS) \\
    Игра «Pong» на стенде SDK-1.1M
    \vfill
    \vspace{\fill}
    \begin{flushright}
        \textbf{Выполнили:} \\
        Козаченко Данил Александрович \\
        Марьин Григорий Алексеевич \\
        \textbf{Группа P3312} \\
        \textbf{Проверил:} \\
        Гончаров Алексей Андреевич
    \end{flushright}
    \vspace{2cm}
    Санкт-Петербург
    2026г.
\end{titlepage}

\setcounter{page}{1}
\begin{centering}
    \renewcommand{\contentsname}{Содержание}
    \tableofcontents
\end{centering}
\newpage

\section{Цель работы}
Реализовать в \code{main.c} игру «Pong» для двух игроков: обработку
клавиш, движение ракеток и мяча, подсчёт очков и отображение состояния
матча. Разделить опрос клавиатуры, игровую логику и вывод на экран
между задачами FreeRTOS.

\section{Задание}
Реализовать на стенде SDK-1.1M игру с двумя ракетками и мячом.
Ракетки располагаются у левой и правой границ поля. Игроки перемещают их
вверх и вниз, отбивая мяч. Пропуск мяча приносит очко противнику;
матч заканчивается при достижении семи очков.

Программа должна выводить легенду перед стартом, поддерживать одновременное
управление двумя игроками, паузу и выход в меню. Дополнительно реализованы
обратный отсчёт перед подачей, ускорение мяча после отбитий, сохранение
лучшей серии отбитий в пределах сеанса работы и звуковое сопровождение победы.
Использование FreeRTOS является обязательной частью реализации.

\begin{table}[H]
\centering\small\caption{Назначение клавиш}\label{tab:keys}
\begin{tabularx}{\linewidth}{@{}>{\centering\arraybackslash}p{22mm}Y@{}}
\toprule
Клавиша & Действие\\\midrule
1 / 7 & Ракетка первого игрока: вверх / вниз\\
3 / 9 & Ракетка второго игрока: вверх / вниз\\
2 & Старт из меню; новый матч после завершения\\
0 & Пауза и продолжение во время игры\\
5 & Выход в меню; прекращение звукового сопровождения\\
\bottomrule
\end{tabularx}
\end{table}

\section{Внесённые изменения в main.c}
В пользовательскую часть \code{main.c} добавлены структура игрового
состояния, функции \code{Pong_*} и обработка действий двух игроков.
Реализованы следующие изменения:

\begin{enumerate}
\item Добавлены структура \code{Pong} и функции запуска матча,
перемещения ракеток, движения мяча, обработки столкновений и подсчёта очков.
\item Созданы задачи \code{KeyboardTask}, \code{GameTask},
\code{DisplayTask} и \code{LogTask}; их запуск вынесен в
\code{Pong_InitRTOS()}.
\item Добавлены очереди \code{input_queue}, \code{frame_queue},
\code{log_queue} и мьютекс \code{i2c_mutex} для взаимодействия задач игры.
\item В \code{Pong_ReadKeys()} реализована маска нажатых клавиш.
Добавлены независимое управление двумя ракетками, обработка удержания
и выделение однократных команд старта, паузы и выхода.
\item Реализована \code{Pong_Draw()} с легендой, игровым полем,
счётом, обратным отсчётом, паузой и экраном победителя.
\item Добавлено воспроизведение первых 15 секунд аудиозаписи после
завершения матча. Фрагмент сжат до 60000 байт; выход в меню и новый
матч останавливают звук.
\end{enumerate}

\subsection{Запуск добавленной логики}
В \code{main()} добавлены вызовы инициализации звука и запуска задач игры.
Функция \code{Pong_InitRTOS()} создаёт очереди, мьютекс и четыре задачи,
после чего передаёт выполнение планировщику.

\codecaption{Добавленные вызовы в main}
\begin{minted}{c}
Audio_Init();
Pong_InitRTOS();
\end{minted}

Периодический опрос и расчёт игрового состояния выполняются внутри
добавленных задач. Обычный цикл \code{while (1)} больше не вызывает
\code{Calculator_Poll()}.

\section{Организация задач и обмена данными}
\subsection{Задачи FreeRTOS}
Для игры добавлены четыре задачи. В таблице~\ref{tab:tasks} приведены
их периоды, приоритеты и выделенные размеры стеков. Опрос клавиатуры
получил приоритет 3, игровая логика --- 2, отображение и журналирование --- 1.

\begin{table}[H]
\centering\small\caption{Параметры прикладных задач}\label{tab:tasks}
\small
\begin{tabularx}{\linewidth}{@{}p{28mm}c c c Y@{}}
\toprule
Задача & Приоритет & Период & Стек, байт & Назначение\\\midrule
\code{KeyboardTask} & 3 & 10 мс & 1536 & Сканирование и подавление дребезга\\
\code{GameTask} & 2 & 10 мс & 2048 & Состояние матча, движение, счёт\\
\code{DisplayTask} & 1 & 40 мс & 2048 & Формирование и передача кадра\\
\code{LogTask} & 1 & По событию & 1024 & Вывод сообщений в USART6\\
\bottomrule
\end{tabularx}
\end{table}

Команда \code{vTaskDelayUntil()} сохраняет плановый период,
отсчитывая его от предыдущего момента запуска. Заданы 100 Гц для ввода
и игровой логики и 25 кадров/с для OLED. Фактическая частота зависит
от времени выполнения операций.

\subsection{Очереди и владение состоянием}
Только \code{GameTask} изменяет структуру \code{Pong}. Другие задачи
передают или получают копии данных, что исключает одновременное изменение
положения мяча и ракеток несколькими задачами.

\begin{table}[H]
\centering\small\caption{Объекты межзадачного взаимодействия}
\begin{tabularx}{\linewidth}{@{}p{29mm}p{30mm}Y@{}}
\toprule
Объект & Размер & Использование\\\midrule
\code{input_queue} & 1 \code{InputState} & Последняя маска клавиш и время опроса\\
\code{frame_queue} & 1 \code{Pong} & Снимок состояния для отрисовки\\
\code{log_queue} & 8 × 64 байта & Сообщения о режимах и забитых мячах\\
\code{i2c_mutex} & Мьютекс & Доступ к общей шине I\textsuperscript{2}C\\
\bottomrule
\end{tabularx}
\end{table}

В \code{input_queue} и \code{frame_queue} записывается последнее состояние.
Запись выполнена через \code{xQueueOverwrite()}, чтение --- через
\code{xQueuePeek()}. Таким образом, ракетки получают актуальную маску
удержания, а экран --- снимок текущего матча.

\begin{figure}[H]
\centering
\begin{tikzpicture}
\node[box] (kb) at (0,0) {\code{KeyboardTask}\\опрос клавиатуры};
\node[box] (game) at (5,0) {\code{GameTask}\\игровое состояние};
\node[box] (disp) at (10,0) {\code{DisplayTask}\\кадр OLED};
\node[box] (log) at (5,-1.8) {\code{LogTask}\\USART6};
\draw[arr] (kb) -- node[above,font=\tiny]{\code{input_queue}} (game);
\draw[arr] (game) -- node[above,font=\tiny]{\code{frame_queue}} (disp);
\draw[arr] (game) -- node[right,font=\tiny]{\code{log_queue}} (log);
\end{tikzpicture}
\caption{Передача копий данных между задачами}
\end{figure}

\clearpage
\section{Диаграмма последовательности}
На рисунке~\ref{fig:sequence} показан один цикл обмена. В реальной программе
периоды задач различаются: несколько шагов игры могут приходиться
на одно обновление экрана. Горизонтальные стрелки не означают прямого
вызова одной задачи другой: состояния передаются через очереди.

\begin{figure}[H]
\centering
\begin{tikzpicture}[x=1cm,y=1cm,font=\small]
\foreach \x/\name in {0/Клавиатура,4/KeyboardTask,8/GameTask,12/DisplayTask}{
 \node[draw,minimum width=29mm,minimum height=9mm,font=\small\ttfamily] at (\x,0) {\name};
 \draw[dashed,black!55] (\x,-0.5)--(\x,-10.8);
}
\draw[arr] (4,-1.0)--node[above]{Выбор строки PCA9538} (0,-1.0);
\draw[arr] (0,-1.8)--node[above]{Состояние столбцов} (4,-1.8);
\node[draw,fill=white,text width=29mm,align=center,font=\footnotesize] at (4,-2.9)
 {Четыре строки;\\дребезг 20 мс};
\draw[arr] (4,-4.0)--node[above]{\code{xQueueOverwrite(input_queue)}} (8,-4.0);
\node[draw,fill=white,text width=30mm,align=center,font=\footnotesize] at (8,-5.1)
 {Ракетки, мяч,\\столкновения, счёт};
\draw[arr] (8,-6.2)--node[above]{\code{xQueueOverwrite(frame_queue)}} (12,-6.2);
\node[draw,fill=white,text width=31mm,align=center,font=\footnotesize] at (12,-7.3)
 {\code{Pong_Draw()}\\поле и счёт};
\draw[arr] (12,-8.4)--++(0.7,0)--++(0,-0.7)--node[below,font=\footnotesize]
 {Передача 8 страниц OLED} (12,-9.1);
\node[draw,fill=white,text width=30mm,align=center,font=\footnotesize] at (8,-10.0)
 {При победе:\\\code{Audio_Start()}};
\end{tikzpicture}
\caption{Опрос клавиатуры, обработка игрового шага и вывод кадра}
\label{fig:sequence}
\end{figure}

Мьютекс охватывает полный цикл выбора и чтения строк клавиатуры.
При обновлении дисплея он захватывается отдельно для каждой страницы.
Таким образом, задача экрана не удерживает шину на время передачи всего
кадра, а клавиатура может получить доступ между страницами.

\section{Обработка клавиатуры и общей шины}
\subsection{Сканирование и одновременные нажатия}
Добавленная функция \code{Pong_ReadKeys()} возвращает маску нажатий
вместо одного символа. Это позволяет одновременно удерживать по одной
клавише каждого игрока. В \code{KeyboardTask} результат передаётся
игровой задаче вместе со временем опроса.

Каждая клавиша представлена битом в 12-разрядной маске:
\begin{center}
\code{bit = 3 * row + column; mask |= 1u << bit;}
\end{center}
Порядок соответствует раскладке
\verb|1 2 3 / 4 5 6 / 7 8 9 / * 0 #|.
Для управления используются биты 0 и 6 у первого игрока,
2 и 8 у второго. Удержание одной клавиши каждым игроком обрабатывается
одновременно. Если у одного игрока нажаты оба направления, их действия
взаимно компенсируются.

Для защиты от возможных фантомных нажатий при обнаружении более двух
клавиш функция возвращает нулевую маску. Это ограничение относится и к попытке нажать командную клавишу
при двух уже удерживаемых направлениях.

\subsection{Подавление дребезга и события команд}
Для каждого бита хранятся предполагаемое состояние и момент его изменения.
Изменение принимается после 20 мс устойчивого состояния. В очередь
передаются подтверждённая маска и время последнего опроса.
Если снимок старше 100 мс, игровая задача считает все кнопки отпущенными,
чтобы ракетка не продолжала двигаться по устаревшим данным.

Перемещение ракеток определяется удержанием, а команды запуска, паузы и
выхода выполняются только по фронту нажатия:
\codecaption{Выделение новых нажатий в GameTask}
\begin{minted}{c}
uint16_t pressed = input.held & ~previous;
previous = input.held;
GameMode old_mode = game.mode;
uint8_t goal = Pong_Tick(&game, input.held, pressed, now);
Audio_OnModeChange(old_mode, game.mode);
\end{minted}

\subsection{Согласование опроса и вывода}
В \code{KeyboardTask} добавлен захват \code{i2c_mutex} на время опроса
матрицы. В \code{DisplayTask} мьютекс захватывается отдельно для каждой
страницы экрана и освобождается после её передачи. Для этого добавлена
функция \code{oled_UpdatePage()}.

Такое разделение позволяет клавиатуре получить доступ к шине между
страницами кадра. Игровая задача при этом продолжает обрабатывать своё
состояние, не ожидая завершения обновления всего экрана.

\clearpage
\section{Игровой алгоритм}
\subsection{Состояния матча}
Структура \code{Pong} хранит режим, координаты ракеток и мяча, компоненты
скорости, счёт, победителя, время подачи и длину серии отбитий.
Переходы между основными состояниями показаны на рисунке~\ref{fig:states}.

\begin{figure}[H]
\centering
\begin{tikzpicture}[x=1cm,y=1cm]
\node[state] (menu) at (0,0) {\code{MENU}\\легенда};
\node[state] (count) at (4.8,0) {\code{COUNTDOWN}\\подача};
\node[state] (play) at (9.6,0) {\code{PLAYING}\\матч};
\node[state] (pause) at (9.6,-2.5) {\code{PAUSED}\\пауза};
\node[state] (end) at (4.8,-4.0) {\code{FINISHED}\\победитель};
\draw[arr] (menu)--node[above,font=\small]{2} (count);
\draw[arr] (count)--node[above,font=\small]{Таймер} (play);
\draw[arr] (play) edge[bend left=22] node[right,font=\small]{0} (pause);
\draw[arr] (pause) edge[bend left=22] node[left,font=\small]{0} (play);
\draw[arr] (play.south west) to[bend left=12] node[below left,font=\small,align=center]
 {Счёт 7} (end.north east);
\draw[arr] (play.north) to[bend right=45] node[above,font=\small]{Гол, счёт меньше 7} (count.north);
\draw[arr] (end.north) -- node[left,font=\small]{2} (count.south);
\draw[arr] (end.west) -| node[below,pos=0.3,font=\small]{5} (menu.south);
\end{tikzpicture}
\caption{Основные переходы; клавиша 5 возвращает в MENU из любого режима}
\label{fig:states}
\end{figure}

Перед первой подачей выдерживается 3 с, после гола --- 1,5 с.
Ожидание реализовано сравнением текущего времени с \code{deadline},
поэтому задачи продолжают работать. На паузе координаты не изменяются.
После победы показываются итоговый счёт и лучшая серия, запускается
15-секундный аудиофрагмент. Новый матч сбрасывает счёт; рекорд серии
сохраняется до перезапуска программы.

\subsection{Движение и столкновения}
Вычисления выполняются с фиксированной точкой: одному пикселю соответствует
\code{Q=512} единиц. При периоде 0,01 с координаты обновляются так:
\begin{center}
\code{x += vx; y += vy;}
\end{center}
Компоненты \code{vx}, \code{vy} выражены в единицах фиксированной точки за шаг,
а не в пикселях в секунду. Начальная горизонтальная скорость равна
230 / (512 × 0,01) ≈ 44,9 пикселя/с. После каждого отбития
её модуль увеличивается на 22 единицы, но ограничивается значением 650
(около 127 пикселей/с). Ракетки перемещаются на 840 единиц за шаг,
то есть примерно 164 пикселя/с.

При ударе о верхнюю или нижнюю границу вертикальная скорость меняет знак.
Столкновение с ракеткой определяется по пересечению её плоскости за шаг.
Высота пересечения вычисляется интерполяцией между старой и новой
позициями, поэтому мяч не пропускает ракетку при увеличении скорости.
Вертикальная скорость после отбития зависит от расстояния до центра
ракетки; почти горизонтальные траектории получают минимальную
вертикальную составляющую. Выход за левую или правую границу даёт очко
противнику. Фрагмент расчёта отбития приведён в приложении~\ref{app:core}.

\section{Добавленные экраны и управление}
В функции \code{Pong_Draw()} добавлены отдельные варианты отрисовки
для меню, подачи, игры, паузы и окончания матча. На поле выводятся
две ракетки, мяч, горизонтальные границы, пунктирная средняя линия
и счёт игроков. Размер ракетки выбран 3 × 14, мяча --- 3 × 3
пикселя. В меню отображается легенда, а после победы --- итоговый счёт
и лучшая серия отбитий.

\begin{figure}[H]
\centering
\begin{tikzpicture}
\begin{scope}[x=0.048cm,y=0.048cm]
 \fill[black] (0,0) rectangle (128,64);
 \node[text=white,inner sep=0pt,anchor=north west,font=\ttfamily\small] at (3,62) {PONG - FIRST TO 7};
 \node[text=white,inner sep=0pt,anchor=north west,font=\ttfamily\small] at (3,51) {P1 1:UP 7:DOWN};
 \node[text=white,inner sep=0pt,anchor=north west,font=\ttfamily\small] at (3,40) {P2 3:UP 9:DOWN};
 \node[text=white,inner sep=0pt,anchor=north west,font=\ttfamily\small] at (3,29) {2:START 0:PAUSE};
 \node[text=white,inner sep=0pt,anchor=north west,font=\ttfamily\small] at (3,18) {5:EXIT  BEST:0};
 \node[text=white,inner sep=0pt,anchor=north west,font=\ttfamily\small] at (3,8) {HOLD KEYS TO MOVE};
\end{scope}
\begin{scope}[xshift=8cm,x=0.048cm,y=0.048cm]
 \fill[black] (0,0) rectangle (128,64);
 \node[text=white,inner sep=0pt,anchor=north west,font=\ttfamily\small] at (2,63) {P1 2};
 \node[text=white,inner sep=0pt,anchor=north,font=\ttfamily\small] at (64,63) {PONG};
 \node[text=white,inner sep=0pt,anchor=north east,font=\ttfamily\small] at (126,63) {1 P2};
 \draw[white] (0,51)--(127,51) (0,1)--(127,1);
 \draw[white,dashed] (64,3)--(64,50);
 \fill[white] (3,24) rectangle (6,38);
 \fill[white] (122,12) rectangle (125,26);
 \fill[white] (79,29) rectangle (82,32);
\end{scope}
\node[font=\small] at (3.07,-0.5) {а) Главное меню};
\node[font=\small] at (11.07,-0.5) {б) Пример игрового поля};
\end{tikzpicture}
\caption{Схематическое представление экранов, составленное по коду отрисовки}
\label{fig:screens}
\end{figure}

Передача кадра разделена на восемь страниц по 128 байт.
\code{DisplayTask} захватывает шину для одной страницы, отправляет её,
освобождает мьютекс и переходит к следующей. Состояние для кадра является
локальной копией: изменения игры во время передачи не меняют уже
сформированный буфер.

\subsection{Порядок использования}
\begin{enumerate}
\item После включения прочитать легенду. Клавиша 2 запускает новый матч.
\item Во время обратного отсчёта можно установить ракетки в начальное положение.
\item Удерживать 1/7 для левой ракетки и 3/9 для правой.
Нажатия разных игроков допускаются одновременно.
\item Клавиша 0 останавливает движение; повторное нажатие продолжает матч.
\item При семи очках появляется экран победителя. Нажать 2 для нового
матча или 5 для возврата к легенде; оба действия останавливают музыку.
\end{enumerate}

Счёт и рекорд не записываются в энергонезависимую память. Возврат в меню
не завершает работу микроконтроллера: программа продолжает опрашивать
клавиатуру и обновлять экран.

\clearpage
\section{Звуковое сопровождение окончания матча}
\subsection{Подготовка и хранение фрагмента}
Для экрана победы использованы первые 15 с предоставленной записи
Morgenshtern «Новый Мерин». Скрипт \code{tools/prepare_audio.py} декодирует
фрагмент, смешивает каналы в моно, фильтрует высокие частоты перед
уменьшением частоты дискретизации и кодирует результат в IMA ADPCM.
По краям фрагмента предусмотрено короткое плавное изменение амплитуды.
Исходный MP3 не хранится в прошивке.

Каждый отсчёт представлен четырьмя битами; в байте находятся два кода,
сначала младший, затем старший. Начальные значения предиктора и индекса
адаптивного шага равны нулю. Размер аудиоданных:
\begin{center}
15 × 8000 = 120000 сэмплов;\\
120000 × 4 / 8 = 60000 байт.
\end{center}
Это примерно 58,6 КиБ; для несжатого 16-разрядного моносигнала той же
длительности потребовалось бы 240000 байт. Массив \code{audio_clip}
объявлен как \code{const} и размещается во Flash, без копирования
всей записи в оперативную память.

\subsection{Таймеры и декодирование}
TIM2 формирует ШИМ на PA5, подключённом к излучателю звука.
TIM6 вызывает обработчик выборок с частотой 8 кГц.
При тактовой частоте таймеров APB1 84 МГц и \code{PSC=0}:
\begin{center}
Частота ШИМ: 84000000 / (2099 + 1) = 40000 Гц;\\
частота выборок: 84000000 / (10499 + 1) = 8000 Гц.
\end{center}
В прерывании выбирается очередной полубайт, восстанавливается знаковый
16-разрядный отсчёт \code{pcm} и изменяется коэффициент заполнения:
\begin{center}
\code{CCR1 = ((pcm + 32768) * 2100) >> 16;}
\end{center}
Декодер хранит только предыдущее значение сигнала и индекс таблицы шага.
Полный PCM-буфер не создаётся. Прерывание имеет приоритет 6, не вызывает
функции FreeRTOS и не выполняет операции I\textsuperscript{2}C.

\subsection{Условия запуска и остановки}
\code{Audio_OnModeChange()} запускает звук только при переходе в
\code{FINISHED}. Повторные шаги в том же состоянии не перезапускают запись.
После обработки 120000 отсчётов обработчик отключает таймер выборок и
звуковой выход. Выход в меню и запуск нового матча прекращают запись
досрочно. Изменение состояния проигрывателя из задачи защищено
критической секцией. Качество ограничено форматом 8 кГц и характеристиками
штатного излучателя; он используется как дополнительное сопровождение игры.

\clearpage
\section{Вывод}
В пользовательской части \code{main.c} реализована игра «Pong» для двух
игроков. Добавлены движение ракеток и мяча, обработка столкновений,
подсчёт очков до семи, обратный отсчёт перед подачей, пауза и выход в меню.
Управление назначено на клавиши 1/7 и 3/9; перед стартом выводится легенда.

Опрос клавиш, игровая логика, формирование кадра и вывод сообщений
распределены между добавленными задачами. Для обмена состоянием
использованы очереди, для согласования опроса и вывода --- мьютекс.

После завершения матча проигрывается 15-секундный аудиофрагмент.
Для уменьшения объёма записи добавлены подготовка моносигнала 8 кГц,
кодирование IMA ADPCM и декодирование при воспроизведении.
Размер аудиоданных составляет 60000 байт. Запуск нового матча и возврат
в меню прекращают воспроизведение.

\clearpage
\appendix
\renewcommand{\reportcodesize}{\fontsize{9}{9.8}\selectfont}
\section{Основные фрагменты программы}\label{app:core}
Фрагменты соответствуют пользовательской части \code{main.c}.
Таблица аудиоданных в приложение не включена: её содержимое генерируется
скриптом подготовки записи и не меняет алгоритм проигрывания.

\subsection{Создание объектов и запуск задач}
\codecaption{Функция Pong\_InitRTOS}
\begin{minted}{c}
static void Pong_InitRTOS(void)
{
  input_queue = xQueueCreate(1, sizeof(InputState));
  frame_queue = xQueueCreate(1, sizeof(Pong));
  log_queue = xQueueCreate(8, 64);
  i2c_mutex = xSemaphoreCreateMutex();
  if (!input_queue || !frame_queue || !log_queue || !i2c_mutex)
    Error_Handler();
  if (xTaskCreate(KeyboardTask, "Keyboard", 384, NULL, 3,
                  &keyboard_task) != pdPASS ||
      xTaskCreate(GameTask, "Pong", 512, NULL, 2,
                  &game_task) != pdPASS ||
      xTaskCreate(DisplayTask, "OLED", 512, NULL, 1,
                  &display_task) != pdPASS ||
      xTaskCreate(LogTask, "UART", 256, NULL, 1, NULL) != pdPASS)
    Error_Handler();
  Pong_Log("PONG FreeRTOS: 1/7=P1 3/9=P2 2=start 0=pause 5=menu");
  vTaskStartScheduler();
  Error_Handler();
}
\end{minted}

\subsection{Обработка команд и игрового шага}
\codecaption{Основные переходы Pong\_Tick}
\begin{minted}{c}
static uint8_t Pong_Tick(Pong *g, uint16_t held,
                         uint16_t pressed, uint32_t now)
{
  if (pressed & KEY_5) { g->mode = MENU; return 0; }
  if (g->mode == MENU || g->mode == FINISHED) {
    if (pressed & KEY_2) Pong_Start(g, now);
    return 0;
  }
  if (pressed & KEY_0) {
    if (g->mode == PLAYING) g->mode = PAUSED;
    else if (g->mode == PAUSED) g->mode = PLAYING;
    return 0;
  }
  if (g->mode == PAUSED) return 0;
  Pong_MovePaddles(g, held);
  if (g->mode == COUNTDOWN) {
    if ((int32_t)(now - g->deadline) >= 0) {
      g->mode = PLAYING; g->countdown = 0;
    } else {
      g->countdown = (uint8_t)((g->deadline - now + 999u) / 1000u);
    }
    return 0;
  }
  return Pong_Physics(g, now);
}
\end{minted}

\clearpage
\subsection{Расчёт отбития мяча}
\codecaption{Изменение скорости при попадании в ракетку}
\begin{minted}{c}
static void Pong_Bounce(Pong *g, int32_t paddle_y, int direction)
{
  int32_t speed = g->vx < 0 ? -g->vx : g->vx;
  int32_t offset = g->ball_y - (paddle_y + (PADDLE_HEIGHT / 2) * Q);
  speed = Clamp(speed + 22, 230, 650);
  g->vx = direction * speed;
  g->vy = Clamp(offset * speed / (8 * Q), -500, 500);
  if (g->vy > -65 && g->vy < 65) g->vy = offset < 0 ? -65 : 65;
  ++g->rally;
  if (g->rally > g->best_rally) g->best_rally = g->rally;
}
\end{minted}

\subsection{Вывод аудио и связь с завершением матча}
\codecaption{Обработчик таймера выборок}
\begin{minted}{c}
void TIM6_DAC_IRQHandler(void)
{
  if (!(TIM6->SR & TIM_SR_UIF)) return;
  TIM6->SR = 0;
  if (!audio_playing) return;
  uint32_t position = audio_sample_pos;
  if (position >= AUDIO_SAMPLE_COUNT) {
    Audio_Disable();
    ++audio_completed;
    return;
  }
  uint8_t packed = audio_clip[position / 2];
  uint8_t code = (position & 1u) ? packed >> 4 : packed & 15u;
  int16_t pcm = ADPCM_Decode(&audio_decoder, code);
  TIM2->CCR1 = ((uint32_t)((int32_t)pcm + 32768)
                * audio_pwm_period) >> 16;
  audio_sample_pos = position + 1;
}
\end{minted}

\codecaption{Однократный запуск и досрочное прекращение записи}
\begin{minted}{c}
static void Audio_OnModeChange(GameMode before, GameMode after)
{
  if (before == after) return;
  if (after == FINISHED) Audio_Start();
  else if (before == FINISHED || after == MENU) Audio_Stop();
}
\end{minted}

\end{document}
